%% =====================================================================
%%  Appendices
%%  Details that would interrupt the main text: proofs, additional results,
%%  source code, and so on. Each \chapter after \appendix becomes
%%  "Appendix A", "Appendix B", ...
%% =====================================================================
\chapter{Source Code}
\label{chap:source-code}

%% Put source code in the lstlisting environment; language= enables syntax colouring.
%% You can also read a file: \lstinputlisting[language=Python]{code/main.py}
\Cref{lst:greedy} shows an implementation of the greedy algorithm (\cref{alg:greedy}) in Python.

\begin{lstlisting}[language=Python, caption={Implementation of the greedy algorithm}, label=lst:greedy]
def greedy_knapsack(values, weights, capacity):
    """Select items in decreasing order of value per weight."""
    order = sorted(range(len(values)),
                   key=lambda i: values[i] / weights[i], reverse=True)
    selected, total_weight = [], 0
    for i in order:
        if total_weight + weights[i] <= capacity:  # select if it fits
            selected.append(i)
            total_weight += weights[i]
    return selected
\end{lstlisting}

\chapter{Notation}
\label{chap:notation}

\Cref{tab:notation} lists the main symbols used in this thesis.
Scalars are written in italic (e.g., $x$), vectors in bold lowercase letters (e.g., $\mathbf{x}$),
matrices in bold capital letters (e.g., $\mathbf{A}$) and sets in calligraphic letters (e.g., $\mathcal{I}$).
When random variables are used, they are written in capital letters (e.g., $X$)
and their values in lowercase letters (e.g., $x$).

%% Follow the notation of the standard textbook in your field
%% (for reinforcement learning, the textbook by Sutton and Barto).
\begin{table}[htbp]
  \centering
  \caption{Main symbols}
  \label{tab:notation}
  \begin{tabular}{ll}
    \toprule
    Symbol & Meaning \\
    \midrule
    $n$ & number of items \\
    $v_i$, $w_i$ & value and weight of item $i$ \\
    $W$ & capacity \\
    $x_i$ & decision variable that equals 1 if item $i$ is selected and 0 otherwise \\
    $\mathbf{x}$ & vector $(x_1, \dots, x_n)$ representing a solution \\
    $f(\mathbf{x})$ & evaluation function of the proposed method \\
    $\lambda$ & penalty weight \\
    \bottomrule
  \end{tabular}
\end{table}
